\documentclass[11pt,a4paper]{article}
\usepackage[margin=2.2cm]{geometry}
\usepackage{amsmath,amssymb}
\usepackage{graphicx}
\usepackage{booktabs}
\usepackage{xcolor}
\usepackage[colorlinks=true,linkcolor=blue!50!black,urlcolor=blue!50!black]{hyperref}
\usepackage{enumitem}
\usepackage{fontspec}
\setmonofont{Menlo}[Scale=0.82]
\setlist{itemsep=2pt,topsep=4pt}
\setlength{\emergencystretch}{3em}

\newcommand{\code}[1]{\texttt{#1}}

\title{Thinning neck test: sharp-interface methods vs.\ standard FD\\[2pt]
\large What happens when a stiff layer necks down to the grid scale}
\author{StagFDTools, branch \texttt{iim-discontinuous-viscosity}}
\date{\today}

\begin{document}
\maketitle

\noindent Boudinage and necking produce stiff layers that thin until they are thinner than a grid
cell. Standard FD with marker-in-cell viscosities handles this smoothly: the layer fades into an
averaged viscosity and can detach. This note tests how the sharp-interface methods of the
branch (IIM stress form order 1 and 2, xGFM; see \code{docs/IIM\_notes},
\code{docs/xGFM\_notes}) behave when the neck width $w$ drops to a few cells and below.

\section{Setup}
\begin{itemize}
  \item Unit box $[-0.5,0.5]^2$, Dirichlet pure shear with layer-parallel extension
    ($v_x = x$, $v_y = -y$).
  \item Horizontal layer of thickness $H = 0.2$, viscosity $\eta_l = 100$ or $10$, in a matrix
    $\eta_m = 1$. The layer has a Gaussian neck,
    $h(x) = \tfrac{H}{2} - \tfrac{H-w}{2}\exp(-((x-x_0)/\sigma)^2)$ with $\sigma = 0.1$. The neck
    centre $(x_0,y_0) = (0.0043, 0.0071)$ is offset from the grid lines.
  \item Methods:
    \begin{itemize}
      \item FD with cell phase fractions as markers would give, using arithmetic or harmonic
        averaging (\code{Stokes} module);
      \item IIM1 and IIM2 (\code{SolveIIMStress});
      \item xGFM (\code{SolveXGFM}).
    \end{itemize}
  \item Reference: xGFM at $n = 640$, so the thinnest neck tested spans at least 8 reference
    cells. Reference values are taken on the same side of $\Gamma$.
  \item Metrics:
    \begin{itemize}
      \item $L^\infty$ and mean $|P|$ errors in the neck zone $|x-x_0| \le 0.15$ and away from it.
        The far zone stops at $|x| \le 0.35$ to avoid the wall/interface corners.
      \item $L^\infty(\mathbf{v})$ errors. For scale, $|\mathbf{v}| \le 0.5$ on the boundary.
      \item The neck thinning rate $-(v_y^{\rm top} - v_y^{\rm bot})/w$ at $x_0$, relative to the
        reference. This is the quantity that matters for time stepping: it controls how fast the
        neck closes.
    \end{itemize}
\end{itemize}
Scripts (in \code{examples/\_Stokes\_IIM}):
\begin{itemize}
  \item \code{NeckThinning.jl}: sweeps $w$ at fixed $n$, Figs.~\ref{fig:sweep100}--\ref{fig:pfields}.
  \item \code{NeckThinning\_convergence.jl}: refines $n$ at fixed $w$,
    Figs.~\ref{fig:conv}--\ref{fig:vfields}.
\end{itemize}

\section{Sweep of the neck width}
Figures~\ref{fig:sweep100} and~\ref{fig:sweep10} show errors against $w/\Delta$ at $n = 40$ and
$80$. Table~\ref{tab:sweep} gives the $n = 40$ values at contrast 100.

\begin{table}[h]
\centering\small
\begin{tabular}{l rrrrr @{\hspace{1.5em}} rrrrr}
\toprule
 & \multicolumn{5}{c}{$L^\infty(\mathbf{v})$, neck zone} & \multicolumn{5}{c}{thinning rate / reference}\\
$w/\Delta$ & 6 & 4 & 2 & 1 & 0.5 & 6 & 4 & 2 & 1 & 0.5\\
\midrule
FD arithmetic & 0.017 & 0.028 & 0.049 & 0.069 & 0.098 & 1.00 & 0.98 & 0.90 & 0.79 & 0.62\\
FD harmonic   & 0.037 & 0.048 & 0.12  & 0.24  & 0.23  & 1.08 & 1.26 & 1.78 & 2.84 & 1.82\\
IIM1          & 0.42  & 0.79  & 1.7   & 1.3   & 1.9   & 2.1  & 3.8  & 3.0  & 5.3  & 11\\
IIM2          & 0.17  & 0.38  & 0.66  & 1.2   & 2.3   & 1.6  & 2.8  & 4.0  & 9.3  & 14\\
xGFM          & 0.13  & 0.25  & 0.31  & 0.71  & 12    & 1.5  & 2.1  & 0.62 & 0.16 & $-59$\\
\bottomrule
\end{tabular}
\caption{Contrast 100, $n = 40$ (from \code{NeckThinning\_eta100.csv}).}
\label{tab:sweep}
\end{table}

\begin{itemize}
  \item \textbf{Contrast 100.}
    \begin{itemize}
      \item FD with arithmetic averaging gives the best velocities and thinning rate at every
        width. It underestimates the rate gently as the neck becomes subgrid: the layer
        ``dissolves'' into the average, as expected from marker-in-cell.
      \item The sharp methods give the most accurate pressure down to $w \approx 2\Delta$, but
        their velocities are $O(1)$ wrong already at $w = 4\Delta$. The neck then thins 2--4
        times too fast, and much faster still below $\Delta$. xGFM even changes sign.
      \item The pressure failure stays local to the neck. The velocity failure spreads into the
        far field.
    \end{itemize}
  \item \textbf{Contrast 10.} All methods are within about 30\% on the rate down to
    $w \approx \Delta$, and IIM2 keeps the best pressure.
  \item The velocity errors of the sharp methods scale roughly with the contrast.
\end{itemize}

\begin{figure}[p]
  \centering
  \includegraphics[width=\textwidth]{figures/NeckThinning_eta100.png}
  \caption{Sweep of the neck width, contrast 100. Circles: $n = 40$; triangles: $n = 80$.}
  \label{fig:sweep100}
  \vspace{1em}
  \includegraphics[width=\textwidth]{figures/NeckThinning_eta10.png}
  \caption{Same as Fig.~\ref{fig:sweep100}, contrast 10.}
  \label{fig:sweep10}
\end{figure}

\begin{figure}[p]
  \centering
  \includegraphics[width=\textwidth]{figures/NeckThinning_fields_eta100.png}
  \caption{Pressure around the neck, $n = 40$, contrast 100, for $w = 4, 2, 1, 0.5\,\Delta$
    (rows). Each panel uses the colour range of the reference; black and magenta mark values
    above and below it. IIM2 best reproduces the localised low pressure in the neck down to
    $2\Delta$. FD smears it over the averaged cells, and harmonic averaging loses it below
    $\Delta$. xGFM develops spurious extremes once the neck is subgrid.}
  \label{fig:pfields}
\end{figure}

\section{Grid convergence at fixed neck width}
Here the neck width is fixed at $w = 0.05$ ($H/4$) and the grid is refined from $n = 20$ to
$160$, so $w/\Delta$ goes $1, 2, 4, 8$ (Fig.~\ref{fig:conv}, Table~\ref{tab:conv}). This is the
situation of a run at a given resolution as the neck closes.

\begin{table}[h]
\centering\small
\begin{tabular}{l rrrr @{\hspace{1.2em}} rrrr @{\hspace{1.2em}} rrrr}
\toprule
 & \multicolumn{4}{c}{mean $|P|$ error} & \multicolumn{4}{c}{$L^\infty(\mathbf{v})$, neck} & \multicolumn{4}{c}{rate / reference}\\
$n$ & 20 & 40 & 80 & 160 & 20 & 40 & 80 & 160 & 20 & 40 & 80 & 160\\
\midrule
\multicolumn{13}{l}{\emph{contrast 100}}\\
FD arith. & 5.5 & 3.1 & 1.8 & 0.57 & 0.064 & 0.049 & 0.029 & 0.017 & 0.79 & 0.90 & 0.96 & 0.94\\
FD harm.  & 36  & 16  & 7.4 & 3.6  & 0.27  & 0.12  & 0.045 & 0.025 & 2.86 & 1.78 & 1.15 & 1.10\\
IIM1      & 14  & 4.2 & 1.1 & 1.1  & 1.3   & 1.7   & 0.40  & 0.28  & 7.3  & 3.0  & 3.3  & 1.9\\
IIM2      & 10  & 1.9 & 0.66 & 0.14 & 1.0  & 0.66  & 0.29  & 0.075 & 7.8  & 4.0  & 2.3  & 1.45\\
xGFM      & 6.5 & 7.3 & 1.4 & 1.5  & 1.2   & 0.31  & 0.23  & 0.044 & 5.6  & 0.62 & 2.0  & 1.15\\
\midrule
\multicolumn{13}{l}{\emph{contrast 10}}\\
FD arith. & 0.49 & 0.33 & 0.19 & 0.060 & 0.025 & 0.018 & 0.011 & 0.0065 & 0.85 & 0.93 & 0.99 & 0.97\\
FD harm.  & 1.35 & 0.73 & 0.38 & 0.17 & 0.061 & 0.031 & 0.014 & 0.0063 & 1.52 & 1.28 & 1.14 & 1.05\\
IIM1      & 0.77 & 0.26 & 0.083 & 0.064 & 0.11 & 0.12 & 0.029 & 0.022 & 1.44 & 1.13 & 1.19 & 1.10\\
IIM2      & 0.59 & 0.11 & 0.031 & 0.0095 & 0.060 & 0.046 & 0.018 & 0.0049 & 1.46 & 1.26 & 1.13 & 1.04\\
xGFM      & 0.49 & 0.36 & 0.075 & 0.083 & 0.084 & 0.042 & 0.016 & 0.0084 & 1.40 & 1.03 & 1.07 & 1.01\\
\bottomrule
\end{tabular}
\caption{Fixed neck width $w = 0.05$; $n = 20, 40, 80, 160$ correspond to $w/\Delta = 1, 2, 4, 8$
  (from \code{NeckThinning\_conv.csv}). Reference thinning rates: 2.54 (contrast 100),
  1.95 (contrast 10).}
\label{tab:conv}
\end{table}

\begin{itemize}
  \item \textbf{All methods converge once the neck is resolved}, but the sharp methods enter
    their asymptotic regime late at high contrast.
    \begin{itemize}
      \item At contrast 100, IIM2 still overestimates the thinning rate by 45\% at
        $w = 8\Delta$, and xGFM by 15\%.
      \item FD with arithmetic averaging is within 10\% from $w = 2\Delta$ on, and within 21\% at
        $w = \Delta$.
    \end{itemize}
  \item \textbf{IIM2 has the best pressure}: the mean $|P|$ error falls by a factor of about 70
    from $n = 20$ to $160$ at contrast 100 (FD arithmetic: 10).
    \begin{itemize}
      \item At contrast 10, IIM2's $L^\infty(P)$ in the neck converges, $6.3 \to 1.1$.
      \item FD's does not ($\approx 16$--$22$): the staircase error near the interface does not
        decrease.
    \end{itemize}
  \item \textbf{xGFM and IIM1 are non-monotone} at contrast 100. For example, xGFM's rate goes
    $5.6, 0.62, 2.0, 1.15$, so a run cannot rely on them giving a consistent thinning rate across
    resolutions.
  \item \textbf{The velocity maps (Fig.~\ref{fig:vfields}) show where the error lives.}
    \begin{itemize}
      \item At contrast 100, the sharp methods give $v_y$ that is too negative above the neck and
        too positive below it, so the neck closes too fast.
      \item Their error is largest ($0.1$--$1$) in the soft matrix directly above and below the
        neck, a column about as wide as the neck. This is the soft side receiving the amplified
        stiff-side error through the jump.
      \item FD's error is smooth and 8--30 times smaller in $L^\infty$: 0.034, against 0.28
        (xGFM), 0.56 (IIM2) and 1.1 (IIM1).
      \item At contrast 10 all methods have errors of a few $10^{-2}$ or less.
    \end{itemize}
\end{itemize}

\begin{figure}[p]
  \centering
  \includegraphics[width=\textwidth]{figures/NeckThinning_conv.png}
  \caption{Fixed neck width $w = 0.05$, grid refinement $n = 20$--$160$ (horizontal axis:
    $w/\Delta$). Top: contrast 100. Bottom: contrast 10. Dashed grey: $O(h)$. Rate panels are
    clipped at 6; the IIM1 and IIM2 values at $w = \Delta$ (7.3, 7.8) are off the plot.}
  \label{fig:conv}
\end{figure}

\begin{figure}[p]
  \centering
  \includegraphics[width=\textwidth]{figures/NeckThinning_vfields_eta100.png}\\[0.5em]
  \includegraphics[width=\textwidth]{figures/NeckThinning_vfields_eta10.png}
  \caption{$v_y$ around the neck (top rows) and its error with respect to the reference (bottom
    rows, $\log_{10}|v_y - v_y^{\rm ref}|$), $n = 40$, $w = 0.05 = 2\Delta$. Upper pair:
    contrast 100. Lower pair: contrast 10. In the $v_y$ panels, black and magenta mark values
    beyond the reference range.}
  \label{fig:vfields}
\end{figure}

\section{Why: contrast amplification}
\begin{itemize}
  \item The jump conditions build the traction on one side of $\Gamma$ from the other side's
    one-sided fit. With the $\eta$-weighted interface trace used by the IIM
    ($v_\Gamma = (\eta_-v_- + \eta_+v_+)/(\eta_-+\eta_+)$), that fit is effectively the stiff side's.
  \item An error $\delta$ in the stiff-side gradient therefore enters the soft side as about
    $(\eta_l/\eta_m)\,\delta$.
  \item In a neck a few cells wide, the stiff-side stencil straddles the whole layer and
    $\delta = O(1)$, so the soft-side velocity error is $O(\eta_l/\eta_m)$ times larger. That is
    consistent with the roughly 10 times larger errors at contrast 100 than at 10.
  \item Thin \emph{stiff} features are the worst case, which is exactly the boudinage
    configuration. Thin weak layers should not amplify (not tested).
  \item Building the jump from the soft side instead (experimental switch
    \code{StokesIIM.TRACE\_WEIGHT[] = :ieta}) fixes the neck: at $w = 4\Delta$ the velocity error
    drops from 0.38 to 0.02. But it destroys a well-resolved stiff inclusion: Schmid \& Podladchikov
    $L^\infty(P)$ goes from 0.19 to $4\times10^3$. So it cannot be used globally.
\end{itemize}

\section{Poisson analogue and cut cells}
The same necked layer was run for $\nabla\cdot(\beta\nabla u) = 0$, $\beta_l/\beta_m = 100$
(script \code{CutCell\_Neck.jl} in \code{examples/\_Poisson\_Interface}, Fig.~\ref{fig:poisson}).
\begin{itemize}
  \item The IIM ``z-form'' with Landry weights, the analogue of the Stokes code, fails in the same
    way, so the problem is not specific to Stokes.
  \item A cut-cell finite-volume scheme degrades gracefully, from 0.003 at $6h$ to 0.013 at
    $0.5h$ in $L^\infty(u)$, and is the most accurate method below about $4h$. It has one unknown
    per phase piece and per cut cell, and a series (harmonic) flux across $\Gamma$, which is
    dominated by the soft side.
  \item A Stokes MAC version of the cut cell (\code{examples/\_Stokes\_IIM/CutCell\_Stokes.jl}) is
    exact for a straight layer, but its interface velocity unknowns are under-determined on curved
    interfaces. It is not usable yet.
\end{itemize}

\begin{figure}[h]
  \centering
  \includegraphics[width=\textwidth]{figures/CutCell_Neck_beta100.png}
  \caption{Poisson neck, $\beta_l/\beta_m = 100$. Left: potential drop across the neck. Right:
    $L^\infty(u)$ away from $\Gamma$. The cut cell is shown in red.}
  \label{fig:poisson}
\end{figure}

\section{Conclusions}
\begin{itemize}
  \item For a stiff neck (contrast $\gtrsim 100$) narrower than about $8\Delta$, the sharp methods
    (IIM1, IIM2, xGFM) are not safe for velocities or thinning rates: a time-dependent necking run
    with them would thin the neck too fast.
  \item FD with arithmetic averaging is robust and slightly underestimates the thinning rate once
    the neck is subgrid, which is the ``natural detachment'' behaviour.
  \item IIM2 remains clearly the best for pressure where the neck is resolved.
  \item At contrast 10 all methods are acceptable down to about $\Delta$.
  \item Options:
    \begin{itemize}
      \item A local switch from sharp treatment to averaged FD below a few cells. The threshold
        depends on the contrast, about $8\Delta$ at 100.
      \item A contrast-robust coupling (cut cell, or CutFEM/Nitsche), which works in Poisson but
        not yet in Stokes.
    \end{itemize}
\end{itemize}

\end{document}
